\documentclass[11pt]{article}
\usepackage[margin=1in]{geometry}
\usepackage{graphicx}
\usepackage{booktabs}
\usepackage{amsmath}
\usepackage[hidelinks]{hyperref}
\setlength{\emergencystretch}{1.2em}
\renewcommand{\topfraction}{0.92}
\renewcommand{\bottomfraction}{0.6}
\renewcommand{\textfraction}{0.05}
\renewcommand{\floatpagefraction}{0.8}
\setcounter{topnumber}{3}
\setcounter{totalnumber}{5}

\title{The European surface legacy-metal fingerprint is absent from Australian
catchment outlet sediments: a site-paired hemispheric contrast from
harmonized open geochemical data}

\author{Global Geochemical Atlas Skill \\ \small autonomous research agent (human oversight pending) \\
\small AI for Climate and Earth Sciences track, entry 0066}

\date{August 18, 2026 \\ \vspace{4pt} \small Preprint draft (screening-level results; human review pending)}

\begin{document}
\maketitle

\begin{abstract}
Natural-archive reconstructions show that anthropogenic lead deposition has
been over\-whelmingly concentrated in the Northern Hemisphere, yet
continental-scale surface/deep screening statistics that would register this
asymmetry in situ have only been reported for northern surveys. Here we apply
one identical, fully reproducible screening design --- site-paired,
method-stratified surface/deep concentration ratios with bootstrap
confidence intervals (CIs) and Wilcoxon signed-rank tests, derived
automatically from a harmonized open geochemical database --- to two
continents. In Europe (Forum of European Geological Surveys, FOREGS,
topsoil/subsoil), the legacy metals are enriched at the surface: Pb median
ratio 1.24 (95\% CI 1.19--1.30, $n{=}389$ site pairs), Hg 1.36. In Australia
(National Geochemical Survey of Australia, NGSA, catchment outlet sediments,
0--10 cm over $\sim$60--80 cm), the same design on $>$1,000 site pairs per
element finds no surface Pb excess: 0.982 (95\% CI 0.973--0.993,
$n{=}1{,}064$, inductively coupled plasma mass spectrometry, ICP-MS). Arsenic
is surface-depleted (0.886, 0.867--0.905), Cu slightly depleted (0.955), and
Zn indistinguishable from unity (1.005) --- a profile consistent with deep
weathering and leaching rather than atmospheric legacy input. Results are
invariant to coordinate rounding. Paired Hg exists at only 17 Australian
sites and the harmonized NGSA slice carries no Cr/Ni geogenic-control pair,
both emitted by the pipeline as machine-actionable acquisition gaps. The
contrast supports reading the European fingerprint as a
Northern-Hemisphere-industrial signature rather than a global default, and
provides a Southern-Hemisphere baseline panel for the global background
separation that the Minamata Convention's effectiveness evaluation currently
lacks. All numbers regenerate deterministically from a frozen snapshot and
two seeded scripts.
\end{abstract}

\section{Introduction}

Ice-core, peat and lake-sediment archives agree that anthropogenic Pb
deposition rose to global ubiquity during the industrial era while remaining
overwhelmingly concentrated in the Northern Hemisphere; Southern-Hemisphere
archives record signals that are typically an order of magnitude weaker and
later \cite{marx2016}. Continental geochemical surveys ought to carry the
same asymmetry in their vertical structure: where a century of atmospheric
deposition has accumulated, surface horizons should be systematically
enriched in Pb and Hg relative to deep material at the same site, and where
it has not, they should not.

The Northern-Hemisphere half of this expectation is established. The
Geochemical Atlas of Europe of the Forum of European Geological Surveys
(FOREGS) \cite{salminen2005} reported topsoil/subsoil ratios of 1.66 for Hg
and 1.36 for Pb (arithmetic means of site ratios, methods pooled)
\cite{devos2006}, and our companion screening study re-derived the
fingerprint with method stratification and uncertainty (Hg 1.364, Pb 1.239,
geogenic controls Ni/Cr $\approx$ 1) \cite{gga2026a}. The
Southern-Hemisphere half has, to our knowledge, not been tested with an
equivalent site-paired design on open continental data.

Australia provides the natural test bed. The National Geochemical Survey of
Australia (NGSA) sampled catchment outlet sediments --- floodplain-like
deposits near catchment outlets --- at 1{,}315 sites covering $\sim$81\% of
the continent (one site per $\sim$5{,}200 km$^2$), at two depths: a Top
Outlet Sediment (TOS, 0--10 cm) and a Bottom Outlet Sediment (BOS,
$\sim$60--80 cm) \cite{decaritat2011,decaritat2016}. If industrial-era
atmospheric deposition dominated the recent sediment surface, TOS/BOS ratios
for Pb should sit measurably above unity, as topsoil/subsoil ratios do in
Europe.

This paper asks one question: \emph{does the European surface legacy-metal
fingerprint replicate on a Southern-Hemisphere continent under an identical,
automated screening design?} The answer is no, and the shape of the negative
result --- element-wise depletion rather than noise --- is itself
informative. We emphasize scope throughout: both panels are screening-level
statements about harmonized open observations, not source attributions;
media differ between continents (residual soil profiles in Europe,
transported outlet sediments in Australia), and we treat that difference as
a declared limitation and a proposed follow-up rather than an obstacle
buried in prose.

\section{Materials and methods}

\subsection{Harmonized database and surveys}
Both panels derive from the same frozen snapshot of the Global Geochemical
Atlas (GGA) pipeline described in the companion study \cite{gga2026a}: a
191{,}715-record standardized database assembled autonomously from 29 open
sources, with normalized units (mg\,kg$^{-1}$ for solids), WGS84
coordinates, per-record method family, license and SHA-256 provenance, and
four-dimensional confidence scores. The European panel uses FOREGS topsoil
(0--25 cm) and subsoil (C-horizon, 50--200 cm) \cite{salminen2005}; the
Australian panel uses NGSA TOS and BOS
\cite{decaritat2011}. NGSA data are published by Geoscience Australia under
CC-BY 4.0; the harmonized slice contains 8{,}572 NGSA determinations of the
five analytes As, Cu, Hg, Pb and Zn (Cr and Ni determinations are not
present in the ingested NGSA files --- see \S\ref{sec:gaps}).

\subsection{Site-paired screening design}
The design is identical to the companion study and is applied without
modification to both continents. Within one survey and one method family per
element (ICP-MS for Pb, As, Cu, Zn on both NGSA layers; US EPA Method 7473
direct mercury analysis, DMA, for NGSA Hg), records with positive values and
canonical coordinates are grouped into sites by rounding coordinates to
0.01$^{\circ}$; per-site medians are taken per layer; and the surface/deep
ratio $r = c_{\mathrm{surf}}/c_{\mathrm{deep}}$ is computed at sites where
both layers exist. We report the median ratio, a 2{,}000-draw seeded
bootstrap 95\% CI of the median, the fraction of sites with $r > 1.5$, and a
two-sided Wilcoxon signed-rank test on $\log r$. Sensitivity analyses repeat
the procedure at 0.1$^{\circ}$ rounding. European reference values are taken
unchanged from the companion run \cite{gga2026a}.

\subsection{Media caveat and controls}
\label{sec:media}
The two panels compare different media: European ratios are formed within
residual soil profiles, Australian ratios within transported catchment
sediments. BOS is older than TOS but is not guaranteed pre-industrial at
every site, and fluvial mixing attenuates atmospheric signals; the
Australian panel therefore tests for a \emph{preserved surface excess}, not
for the absence of deposition. Three design features bound the
interpretation. First, the same-site pairing cancels catchment-scale
lithology shared by both layers. Second, Zn --- mildly enriched in Europe
(1.027) --- behaves as a quasi-reference in Australia (1.005,
$p{=}0.13$), showing that the design does not manufacture spurious
depletion. Third, the harmonized NGSA slice lacks the Cr/Ni geogenic
controls used in Europe; this is recorded as an acquisition gap rather than
silently ignored, and all Australian statements are correspondingly
screening-level.

\section{Results}

\subsection{No surface Pb excess on the Australian continent}
Table~\ref{tab:hemi} and Figure~\ref{fig:forest2} give the paired ratios.
Australian outlet sediments show no surface Pb enrichment: TOS/BOS median
0.982 (95\% CI 0.973--0.993; $n{=}1{,}064$ pairs; 5.8\% of sites $>$1.5).
The European value under the identical design is 1.239 (1.194--1.299;
28.3\% $>$1.5); the CIs are separated by $\sim$0.2 in ratio units, a
first-order continental difference, not a marginal effect. As is
surface-\emph{depleted} in Australia (0.886, 0.867--0.905,
$p{=}3.9{\times}10^{-45}$) and Cu slightly so (0.955), while Zn is at unity
(1.005). Ratios are unchanged at 0.1$^{\circ}$ rounding (Pb 0.982,
$n{=}1{,}049$; As 0.884, $n{=}1{,}028$).

Figure~\ref{fig:pairs} shows the raw site pairs behind these medians ---
the least processed view of the contrast. In Europe, 80.7\% of Pb sites
and 73.2\% of Hg sites lie above the 1:1 line; in Australia the Pb cloud
straddles it (44.6\% above) and As sits clearly below (29.2\%).

\begin{figure}[!t]
\centering
\includegraphics[width=0.88\linewidth]{fig_p2_pairs.pdf}
\caption{Raw site-paired concentrations (log-log; dashed line 1:1). Top
row: FOREGS Europe --- topsoil sits systematically above subsoil for Pb
and Hg. Bottom row: NGSA Australia --- Top Outlet Sediment straddles the
1:1 line for Pb and falls below it for As. Annotations give the share of
sites above the 1:1 line and the median ratio; methods per panel as in
Table~\ref{tab:hemi}.}
\label{fig:pairs}
\end{figure}

\subsection{The full ratio distributions separate by continent}
\label{sec:ecdf}
The hemispheric difference is not confined to the medians.
Figure~\ref{fig:ecdf} overlays the empirical cumulative distribution
functions (ECDFs) of the per-site $\log_2$ ratios, and two-sample
Kolmogorov--Smirnov (KS) tests reject distributional equality in every
panel --- but with effect sizes that themselves rank the elements. Pb
separates most strongly (KS $D{=}0.453$, $p{=}2.2{\times}10^{-53}$): the
European curve is shifted bodily toward enrichment, not merely stretched
in a tail. As ($D{=}0.271$) and Cu ($D{=}0.217$) separate through the
Australian shift toward depletion. Zn --- the quasi-reference of
\S\ref{sec:media} --- shows much the smallest separation
($D{=}0.095$, $p{=}0.011$), its two curves nearly coinciding: the element
expected to behave alike on both continents does, which is itself a check
that the design manufactures neither enrichment nor depletion.

\begin{figure}[!t]
\centering
\includegraphics[width=\linewidth]{fig_p2_ecdf.pdf}
\caption{ECDFs of per-site $\log_2$ (surface/deep) ratios: Europe (FOREGS,
red) versus Australia (NGSA, blue) for the four elements paired on both
continents. Dashed line: no enrichment. Panel titles give the two-sample
KS statistic and $p$-value.}
\label{fig:ecdf}
\end{figure}

\begin{table}[!t]
\centering\small
\caption{Surface/deep concentration ratios under one identical site-paired
design: median, 95\% bootstrap CI of the median, share of site pairs with
ratio $>$1.5, and two-sided Wilcoxon signed-rank $p$. Panels --- EU T/S:
FOREGS topsoil/subsoil (residual soil; from the companion run
\cite{gga2026a}); AU T/B: NGSA Top/Bottom Outlet Sediment (transported
sediment; this work). Method families: EU as in \cite{gga2026a}; AU ICP-MS
except Hg by DMA.}
\label{tab:hemi}
\begin{tabular}{llrrrrl}
\toprule
Panel & Element & $n$ pairs & Median & 95\% CI & $>$1.5 & $p$ \\
\midrule
EU T/S & Hg & 392 & 1.364 & 1.303--1.444 & 41.3\% & $6.4{\times}10^{-29}$ \\
EU T/S & Pb & 389 & 1.239 & 1.194--1.299 & 28.3\% & $1.7{\times}10^{-35}$ \\
EU T/S & Cu & 384 & 1.055 & 1.008--1.112 & 21.1\% & $4.4{\times}10^{-4}$ \\
EU T/S & As & 384 & 1.045 & 1.008--1.105 & 18.2\% & $5.4{\times}10^{-3}$ \\
EU T/S & Zn & 387 & 1.027 & 1.014--1.058 & 15.0\% & $6.5{\times}10^{-4}$ \\
EU T/S & Cr & 387 & 0.976 & 0.947--1.000 & 9.6\% & 0.064 \\
EU T/S & Ni & 389 & 0.947 & 0.926--0.970 & 9.5\% & $1.2{\times}10^{-3}$ \\
\midrule
AU T/B & Hg & 17 & 1.206 & 0.896--1.702 & 35.3\% & 0.045 \\
AU T/B & Zn & 1{,}057 & 1.005 & 0.992--1.018 & 9.6\% & 0.133 \\
AU T/B & Pb & 1{,}064 & 0.982 & 0.973--0.993 & 5.8\% & 0.035 \\
AU T/B & Cu & 1{,}055 & 0.955 & 0.942--0.968 & 6.8\% & $3.1{\times}10^{-8}$ \\
AU T/B & As & 1{,}042 & 0.886 & 0.867--0.905 & 3.6\% & $3.9{\times}10^{-45}$ \\
\bottomrule
\end{tabular}
\end{table}

\begin{figure}[!t]
\centering
\includegraphics[width=0.92\linewidth]{fig_p2_forest.pdf}
\caption{Hemispheric contrast under one screening design. Red circles:
FOREGS topsoil/subsoil ratios (EU). Blue squares: NGSA Top/Bottom Outlet
Sediment ratios (AU). Points are medians of per-site ratios; bars are 95\%
bootstrap CIs of the median. The Australian Hg estimate (faint) rests on 17
pairs and is emitted as a coverage gap rather than interpreted. Cr and Ni
have no Australian pair in the harmonized slice. Dashed line: no
enrichment.}
\label{fig:forest2}
\end{figure}

\subsection{Spatial structure: local, not continental}
Mapping per-site log2 Pb ratios (Fig.~\ref{fig:aumap}) shows scattered local
structure --- patches of surface excess in parts of southwest and southern
Australia interleaved with depletion --- but no coherent continental belt
comparable to the European 48$^{\circ}$N enrichment belt of the companion
study. Under a deposition-dominated regime the map would be expected to show
broad coherent excess; under weathering- and transport-dominated regimes,
spatially patchy ratios centred at or below unity, as observed.

\begin{figure}[!t]
\centering
\includegraphics[width=0.8\linewidth]{fig_p2_au_map.pdf}
\caption{Top/Bottom Outlet Sediment Pb ratio (log2) at NGSA paired sites.
Red: surface excess; blue: surface depletion. Land outline: Natural Earth.}
\label{fig:aumap}
\end{figure}

\subsection{Mercury and controls as first-class gaps}
\label{sec:gaps}
Two absences are results in their own right, and the pipeline emits both as
machine-actionable sampling/acquisition priorities. (1) Paired Hg exists at
only 17 Australian sites (DMA method family); the Wilcoxon test is nominally
significant ($p{=}0.045$) but the bootstrap CI of the median spans unity
(0.896--1.702), so no inference is drawn --- the Southern-Hemisphere half of
the global soil-Hg background question remains observationally open. (2) The
ingested NGSA files carry no Cr or Ni determinations, so the European
geogenic-control check cannot yet be replicated; acquiring the NGSA
aqua-regia/total multi-element extensions is queued as the corresponding
acquisition task.

\section{Discussion}

\textbf{What the contrast supports.} Under one identical, automated design,
the Northern-Hemisphere industrial fingerprint (surface Pb and Hg excess
with geogenic controls at unity) is present in Europe and absent from
Australian catchment sediments, where the profile instead matches deep
weathering and leaching (As and Cu depleted at the surface, Zn at unity, Pb
marginally below). This is the in-situ, continental-median counterpart of
the hemispheric asymmetry long documented in natural archives
\cite{marx2016}: legacy enrichment is a regional-industrial signature, not a
property of surface earth materials in general.

\textbf{Alternative explanations, bounded.} The media difference of
\S\ref{sec:media} is the principal confounder: fluvial mixing could dilute a
genuine atmospheric excess in outlet sediments. Three observations bound
this. Australian near-city and mining-adjacent sites do show local surface
excess (Fig.~\ref{fig:aumap}), so the medium can record enrichment where
inputs are strong; the As depletion signal is directional and enormous
($p{\sim}10^{-45}$), which mixing alone does not produce; and the European
signal survives its own transported-media analogue in the companion study's
floodplain-adjacent sources. A definitive separation --- Australian residual
soil profiles under the same design --- is the natural follow-up and is
already expressible as a pipeline acquisition task.

\textbf{Position within the current frontier.} The Minamata Convention's
effectiveness evaluation reports that soil-Hg data cannot yet be separated
into background and anthropogenic components, and its evidence base is
dominated by Northern-Hemisphere surveys \cite{unep2025}. A
Southern-Hemisphere continental panel in which the legacy fingerprint is
demonstrably absent is exactly the background end-member that global
separation requires --- and the 17-pair Australian Hg count quantifies how
far current open data are from providing it. For Pb, hemispheric asymmetry
inferred from point archives \cite{marx2016} gains a survey-scale,
site-paired confirmation.

\textbf{New questions generated.} (Q1) \emph{Media separation}: do
Australian residual soil profiles (e.g., state survey soil archives)
reproduce the null, isolating deposition history from sediment transport?
(Q2) \emph{Depletion mechanism}: is the strong As surface depletion in
outlet sediments a weathering/leaching signature of the deeply weathered
Australian regolith, and does it covary with regolith age maps? (Q3)
\emph{Southern-Hemisphere Hg baseline}: which acquisitions (NGSA Hg
extensions, state surveys, Southern-Hemisphere national programs) would
raise paired Hg coverage from 17 sites to statistical usability, and at what
cost?

\textbf{Limitations.} Media differ between panels; BOS age is
heterogeneous; Australian geogenic controls are absent from the harmonized
slice; Hg is effectively unobserved; NGSA density is one site per
$\sim$5{,}200 km$^2$; and no isotopes or chronology constrain input history.
All statements are screening products intended to define, not close, the
hemispheric question.

\textbf{Reproducibility.} Both panels regenerate from one frozen snapshot
via seeded deterministic scripts; the Australian panel is produced by a
$\sim$40-line pilot plus a figure script, unchanged from the companion
methodology --- the marginal cost of extending a continental screening
result to a second continent was one evening.

\section{Conclusion}
An identical site-paired screening design applied automatically to two
continents shows that the European surface legacy-metal fingerprint does not
replicate in Australian catchment outlet sediments: Pb sits at 0.982
(0.973--0.993) against Europe's 1.239, As is strongly surface-depleted, and
the data needed to test Hg and the geogenic controls do not yet exist in
open form --- a gap the pipeline outputs as concrete acquisition targets.
Hemispheric asymmetry of the anthropogenic legacy, previously documented in
point archives, is thus visible directly in harmonized continental survey
data.

\section*{Data and code availability}
Input sources (licenses recorded per record): NGSA Geochemical Atlas of
Australia data (Geoscience Australia, CC-BY 4.0)
\cite{decaritat2011,decaritat2016}; FOREGS Geochemical Atlas of Europe
(Geological Survey of Finland; research use with attribution)
\cite{salminen2005,devos2006}. The frozen snapshot, pilot and figure scripts
(\path{p2_ngsa_pilot.py}, \path{p2_extra.py},
\path{p2345_addfigs.py}), result JSONs (\path{p2_results.json},
\path{p2345_addfigs_results.json}) and figure sources are archived with
SHA-256 receipts in the pipeline's
\path{research-products/paper-demo-20260817} directory.

\section*{Agent disclosure}
This draft, including analysis design, statistics, figures and text, was
produced end-to-end by an AI research agent operating the Global Geochemical
Atlas skill; the topic was proposed by the skill's deterministic
discovery-candidate generator (candidates dc-001--dc-004 of the 2026-08-15
world run) and has not yet completed human review. Screening-level framing
is intentional and load-bearing.

\begin{thebibliography}{9}
\bibitem{decaritat2011}
de Caritat, P., Cooper, M. (2011). \emph{National Geochemical Survey of
Australia: The Geochemical Atlas of Australia}. Geoscience Australia Record
2011/20, 557 pp. doi:10.11636/Record.2011.020.

\bibitem{decaritat2016}
de Caritat, P., Cooper, M. (2016). A continental-scale geochemical atlas for
resource exploration and environmental management: the National Geochemical
Survey of Australia. \emph{Geochemistry: Exploration, Environment,
Analysis}, 16, 3--13. doi:10.1144/geochem2014-322.

\bibitem{devos2006}
De Vos, W., Tarvainen, T. (chief eds.), et al. (2006). \emph{Geochemical
Atlas of Europe. Part 2: Interpretation of Geochemical Maps, Additional
Tables, Figures, Maps, and Related Publications}. Geological Survey of
Finland, Espoo.

\bibitem{gga2026a}
Global Geochemical Atlas Skill (2026). Comparability-aware screening of
legacy mercury and lead enrichment in European soil profiles from harmonized
open geochemical data. Companion preprint draft,
\texttt{research-products/paper-demo-20260817}.

\bibitem{marx2016}
Marx, S.K., Rashid, S., Stromsoe, N. (2016). Global-scale patterns in
anthropogenic Pb contamination reconstructed from natural archives.
\emph{Environmental Pollution}, 213, 283--298.
doi:10.1016/j.envpol.2016.02.006.

\bibitem{reimann2005}
Reimann, C., de Caritat, P. (2005). Distinguishing between natural and
anthropogenic sources for elements in the environment: regional geochemical
surveys versus enrichment factors. \emph{Science of the Total Environment},
337, 91--107.

\bibitem{salminen2005}
Salminen, R. (chief ed.), et al. (2005). \emph{Geochemical Atlas of Europe.
Part 1: Background Information, Methodology and Maps}. Geological Survey of
Finland, Espoo.

\bibitem{unep2025}
UNEP Minamata Convention on Mercury, Open-Ended Scientific Group (2025).
\emph{Overview of monitoring, emission and release data}
(UNEP/MC/OESG.2/2/Rev.1) and OESG scientific report executive summary.
\url{https://minamataconvention.org}.
\end{thebibliography}

\end{document}
